\documentclass[11pt]{article}
\usepackage[margin=1in]{geometry}
\usepackage{amsmath,amssymb,amsthm,mathrsfs,microtype}
\usepackage[hidelinks]{hyperref}
\newtheorem{theorem}{Theorem}
\newcommand{\B}{\mathcal B}
\newcommand{\domn}{\operatorname{dom}}
\title{The two remaining domain pathologies in properly infinite von Neumann algebras}
\author{Autonomous proof completion for arXiv:2311.16170}
\date{11 August 2026}

\begin{document}
\maketitle

\begin{abstract}
Ghosh and Nayak ask three domain-pathology questions for affiliated operators
of a properly infinite von Neumann algebra, answer the third, and leave the
first two to the reader.  We give the short complete argument.  Their normal
embedding of $\B(\ell^2)$ and canonical affiliated-operator functor transport
von Neumann's classical disjoint-domain pair and Chernoff's square-trivial
symmetric operator while preserving products, domains, lattice intersections,
closedness, and the relevant order and adjoint properties.
\end{abstract}

\section{Statement}

Let $\mathscr M\subseteq\B(\mathcal H)$ be a properly infinite von Neumann
algebra and let $\mathscr M_{\rm aff}^{c}$ denote its closed affiliated
operators in the notation of the source paper.

\begin{theorem}\label{thm:main}
Both Questions 1 and 2 of Section 8 of the source have affirmative answers:
\begin{enumerate}
\item There are a positive self-adjoint $T\in\mathscr M_{\rm aff}^{c}$ and a
unitary $U\in\mathscr M$ such that
\[
 \domn(T)\cap\domn(UTU^*)=\{0\}_{\mathcal H}.
\]
\item There is a densely-defined closed symmetric
$A\in\mathscr M_{\rm aff}^{c}$ such that
\[
 \domn(A^2)=\{0\}_{\mathcal H}.
\]
\end{enumerate}
\end{theorem}

\section{The transport identities}

By Proposition 8.1 of the source there is a unital normal embedding
\[
 \Phi:\B(\ell^2(\mathbb N))\longrightarrow\mathscr M.
\]
Write $\Phi_{\rm aff}$ for its canonical extension to affiliated operators
and $\Phi_{\rm aff}^s$ for its extension to affiliated operator ranges.  The
source proves the following facts:
\begin{align}
 \domn(\Phi_{\rm aff}(S))
   &=\Phi_{\rm aff}^s(\domn S),                                      \label{eq:domain}\\
 \Phi_{\rm aff}(S_1S_2)
   &=\Phi_{\rm aff}(S_1)\Phi_{\rm aff}(S_2),                         \label{eq:product}\\
 \Phi_{\rm aff}(S^*)&=\Phi_{\rm aff}(S)^*.                          \label{eq:adjoint}
\end{align}
Here \eqref{eq:domain} is Theorem 5.5(i), and
\eqref{eq:product}--\eqref{eq:adjoint} are Theorem 4.14(ii),(iv).
Theorem 3.10 says that $\Phi_{\rm aff}^s$ preserves intersections and
\begin{equation}
\label{eq:zero}
 \Phi_{\rm aff}^s(\{0\}_{\ell^2})=\{0\}_{\mathcal H}.
\end{equation}
Theorem 5.5 preserves dense definition and closedness; Theorem 6.6 preserves
symmetry, positivity, and self-adjointness.

\section{Question 1}

The classical theorem of von Neumann cited by the source supplies an
unbounded positive self-adjoint operator $T_0$ and a unitary $U_0$ on
$\ell^2(\mathbb N)$ such that
\begin{equation}\label{eq:vn}
 \domn(T_0)\cap\domn(U_0T_0U_0^*)=\{0\}_{\ell^2}.
\end{equation}
(This is the reason the source calls Question 1 a weaker version of von
Neumann's result.)  Put
\[
 T=\Phi_{\rm aff}(T_0),\qquad U=\Phi(U_0).
\]
Because $\Phi$ is a unital $*$-homomorphism, $U$ is unitary.  The preservation
results above show that $T$ is positive, self-adjoint, closed, and affiliated,
and repeated use of \eqref{eq:product}--\eqref{eq:adjoint} gives
\[
 UTU^*=\Phi_{\rm aff}(U_0T_0U_0^*).
\]
Consequently, by \eqref{eq:domain}, lattice preservation, \eqref{eq:vn}, and
\eqref{eq:zero},
\begin{align*}
 \domn(T)\cap\domn(UTU^*)
 &=\Phi_{\rm aff}^s(\domn T_0)
     \cap\Phi_{\rm aff}^s(\domn(U_0T_0U_0^*))\\
 &=\Phi_{\rm aff}^s\!\left(
      \domn T_0\cap\domn(U_0T_0U_0^*)\right)
  =\{0\}_{\mathcal H}.
\end{align*}

\section{Question 2}

Chernoff's explicit example cited by the source is a densely-defined,
semibounded, closed symmetric operator $A_0$ on $\ell^2(\mathbb N)$ such that
\begin{equation}\label{eq:chernoff}
 \domn(A_0^2)=\{0\}_{\ell^2}.
\end{equation}
Set $A=\Phi_{\rm aff}(A_0)$.  Theorems 5.5 and 6.6 show that $A$ is
densely-defined, closed, symmetric, and affiliated.
By exact product preservation,
\[
 A^2=\Phi_{\rm aff}(A_0^2).
\]
Using \eqref{eq:domain}, \eqref{eq:chernoff}, and \eqref{eq:zero} now gives
\[
 \domn(A^2)
 =\Phi_{\rm aff}^s(\domn(A_0^2))
 =\Phi_{\rm aff}^s(\{0\}_{\ell^2})
 =\{0\}_{\mathcal H}.
\]
This proves Theorem~\ref{thm:main}. \qed

\section*{Scope and references}

The exact questions occur on source PDF pages 36--37; the source explicitly
leaves them to the reader on page 37.  The Hilbert-space inputs are J. von
Neumann, \emph{Zur Theorie der unbeschr\"ankten Matrizen}, J. Reine Angew.
Math. 161 (1929), 208--236, and P. R. Chernoff, \emph{A semibounded closed
symmetric operator whose square has trivial domain}, Proc. Amer. Math. Soc.
89 (1983), 289--290.  The proof is a completion of the indicated exercise,
not a claim that those classical pathologies are new.

\end{document}
